% Standalone file -- % Standalone file -- % Standalone file -- % Standalone file -- \input{phase3_solver_table} in the Phase 3 report.
% Requires: booktabs, float (already loaded in the Phase 2 main doc)

\begin{table}[htbp]
\centering
\small
\begin{tabular}{llllll}
\toprule
\textbf{Scenario} & \textbf{Method} & \textbf{Siting Buses} & \textbf{24h Cost (\$)} & \textbf{Runtime} & \textbf{Gap vs.\ Classical} \\
\midrule
\multirow{3}{*}{No DC}
  & Classical siting (Benders/SCIP) & (1, 2, 3, 8)   & 97,229  & 16.3s  & (reference) \\
  & Quantum VQA --- Aer statevector ($n=20$) & (8, 11, 13, 14) & 97,229  & 102.4s & 0.000\% \\
  & QPU (IonQ Forte-1)               & ---            & ---     & ---    & --- \\
\addlinespace
\multirow{3}{*}{DC @ bus 4}
  & Classical siting (Benders/SCIP) & (3, 8, 11, 12) & 193,229 & 18.1s  & (reference) \\
  & Quantum VQA --- Aer statevector ($n=20$) & (4, 10, 11, 12) & 193,229 & 139.4s & 0.000\% \\
  & QPU (IonQ Forte-1)               & ---            & ---     & ---    & --- \\
\addlinespace
\multirow{3}{*}{DC @ 4, Gen2 outage}
  & Classical siting (Benders/SCIP) & (1, 2, 6, 7)   & 227,343 & 18.0s  & (reference) \\
  & Quantum VQA --- Aer statevector ($n=20$) & (4, 6, 9, 12) & 227,343 & 83.8s  & 0.000\% \\
  & QPU (IonQ Forte-1)               & ---            & ---     & ---    & --- \\
\addlinespace
\multirow{3}{*}{DC @ 4, heatwave}
  & Classical siting (Benders/SCIP) & (7, 9, 11, 12) & 203,675 & 34.6s  & (reference) \\
  & Quantum VQA --- Aer statevector ($n=20$) & (4, 9, 10, 14) & 203,404 & 76.8s  & $-$0.133\%\textsuperscript{\dag} \\
  & QPU (IonQ Forte-1)               & ---            & ---     & ---    & --- \\
\bottomrule
\end{tabular}
\caption{Cost and runtime comparison across four IEEE14 scenarios. Classical
siting uses Benders decomposition with SCIP subproblems (batch-Benders,
2-min time limit, early-stall detection). Quantum VQA uses a 19-qubit
butterfly ansatz (COBYLA, Aer CPU statevector simulator), evaluating the top
$n=20$ candidates from the quantum sieve classically (UC second stage). Gap
vs.\ Classical is the percentage cost difference relative to the classical
result; 0.000\% means the two methods matched to the nearest cent at this
problem size. QPU rows are placeholders (real Forte-1 hardware runs not yet
executed for this comparison --- budget-limited to a small number of real
QPU submissions). Gen2 outage (full 24h horizon) and the heatwave demand
scaling are independent scenarios, not combined. \textsuperscript{\dag}Negative
gap: the classical Benders search stopped early via its stall-detection
heuristic (no cost improvement for several iterations) before reaching the
placement the quantum sieve's candidate list happened to include, so quantum
found a slightly cheaper feasible placement than the classical run's
reported best for this scenario --- not a discrepancy in either solver's
dispatch model.}
\end{table}
 in the Phase 3 report.
% Requires: booktabs, float (already loaded in the Phase 2 main doc)

\begin{table}[htbp]
\centering
\small
\begin{tabular}{llllll}
\toprule
\textbf{Scenario} & \textbf{Method} & \textbf{Siting Buses} & \textbf{24h Cost (\$)} & \textbf{Runtime} & \textbf{Gap vs.\ Classical} \\
\midrule
\multirow{3}{*}{No DC}
  & Classical siting (Benders/SCIP) & (1, 2, 3, 8)   & 97,229  & 16.3s  & (reference) \\
  & Quantum VQA --- Aer statevector ($n=20$) & (8, 11, 13, 14) & 97,229  & 102.4s & 0.000\% \\
  & QPU (IonQ Forte-1)               & ---            & ---     & ---    & --- \\
\addlinespace
\multirow{3}{*}{DC @ bus 4}
  & Classical siting (Benders/SCIP) & (3, 8, 11, 12) & 193,229 & 18.1s  & (reference) \\
  & Quantum VQA --- Aer statevector ($n=20$) & (4, 10, 11, 12) & 193,229 & 139.4s & 0.000\% \\
  & QPU (IonQ Forte-1)               & ---            & ---     & ---    & --- \\
\addlinespace
\multirow{3}{*}{DC @ 4, Gen2 outage}
  & Classical siting (Benders/SCIP) & (1, 2, 6, 7)   & 227,343 & 18.0s  & (reference) \\
  & Quantum VQA --- Aer statevector ($n=20$) & (4, 6, 9, 12) & 227,343 & 83.8s  & 0.000\% \\
  & QPU (IonQ Forte-1)               & ---            & ---     & ---    & --- \\
\addlinespace
\multirow{3}{*}{DC @ 4, heatwave}
  & Classical siting (Benders/SCIP) & (7, 9, 11, 12) & 203,675 & 34.6s  & (reference) \\
  & Quantum VQA --- Aer statevector ($n=20$) & (4, 9, 10, 14) & 203,404 & 76.8s  & $-$0.133\%\textsuperscript{\dag} \\
  & QPU (IonQ Forte-1)               & ---            & ---     & ---    & --- \\
\bottomrule
\end{tabular}
\caption{Cost and runtime comparison across four IEEE14 scenarios. Classical
siting uses Benders decomposition with SCIP subproblems (batch-Benders,
2-min time limit, early-stall detection). Quantum VQA uses a 19-qubit
butterfly ansatz (COBYLA, Aer CPU statevector simulator), evaluating the top
$n=20$ candidates from the quantum sieve classically (UC second stage). Gap
vs.\ Classical is the percentage cost difference relative to the classical
result; 0.000\% means the two methods matched to the nearest cent at this
problem size. QPU rows are placeholders (real Forte-1 hardware runs not yet
executed for this comparison --- budget-limited to a small number of real
QPU submissions). Gen2 outage (full 24h horizon) and the heatwave demand
scaling are independent scenarios, not combined. \textsuperscript{\dag}Negative
gap: the classical Benders search stopped early via its stall-detection
heuristic (no cost improvement for several iterations) before reaching the
placement the quantum sieve's candidate list happened to include, so quantum
found a slightly cheaper feasible placement than the classical run's
reported best for this scenario --- not a discrepancy in either solver's
dispatch model.}
\end{table}
 in the Phase 3 report.
% Requires: booktabs, float (already loaded in the Phase 2 main doc)

\begin{table}[htbp]
\centering
\small
\begin{tabular}{llllll}
\toprule
\textbf{Scenario} & \textbf{Method} & \textbf{Siting Buses} & \textbf{24h Cost (\$)} & \textbf{Runtime} & \textbf{Gap vs.\ Classical} \\
\midrule
\multirow{3}{*}{No DC}
  & Classical siting (Benders/SCIP) & (1, 2, 3, 8)   & 97,229  & 16.3s  & (reference) \\
  & Quantum VQA --- Aer statevector ($n=20$) & (8, 11, 13, 14) & 97,229  & 102.4s & 0.000\% \\
  & QPU (IonQ Forte-1)               & ---            & ---     & ---    & --- \\
\addlinespace
\multirow{3}{*}{DC @ bus 4}
  & Classical siting (Benders/SCIP) & (3, 8, 11, 12) & 193,229 & 18.1s  & (reference) \\
  & Quantum VQA --- Aer statevector ($n=20$) & (4, 10, 11, 12) & 193,229 & 139.4s & 0.000\% \\
  & QPU (IonQ Forte-1)               & ---            & ---     & ---    & --- \\
\addlinespace
\multirow{3}{*}{DC @ 4, Gen2 outage}
  & Classical siting (Benders/SCIP) & (1, 2, 6, 7)   & 227,343 & 18.0s  & (reference) \\
  & Quantum VQA --- Aer statevector ($n=20$) & (4, 6, 9, 12) & 227,343 & 83.8s  & 0.000\% \\
  & QPU (IonQ Forte-1)               & ---            & ---     & ---    & --- \\
\addlinespace
\multirow{3}{*}{DC @ 4, heatwave}
  & Classical siting (Benders/SCIP) & (7, 9, 11, 12) & 203,675 & 34.6s  & (reference) \\
  & Quantum VQA --- Aer statevector ($n=20$) & (4, 9, 10, 14) & 203,404 & 76.8s  & $-$0.133\%\textsuperscript{\dag} \\
  & QPU (IonQ Forte-1)               & ---            & ---     & ---    & --- \\
\bottomrule
\end{tabular}
\caption{Cost and runtime comparison across four IEEE14 scenarios. Classical
siting uses Benders decomposition with SCIP subproblems (batch-Benders,
2-min time limit, early-stall detection). Quantum VQA uses a 19-qubit
butterfly ansatz (COBYLA, Aer CPU statevector simulator), evaluating the top
$n=20$ candidates from the quantum sieve classically (UC second stage). Gap
vs.\ Classical is the percentage cost difference relative to the classical
result; 0.000\% means the two methods matched to the nearest cent at this
problem size. QPU rows are placeholders (real Forte-1 hardware runs not yet
executed for this comparison --- budget-limited to a small number of real
QPU submissions). Gen2 outage (full 24h horizon) and the heatwave demand
scaling are independent scenarios, not combined. \textsuperscript{\dag}Negative
gap: the classical Benders search stopped early via its stall-detection
heuristic (no cost improvement for several iterations) before reaching the
placement the quantum sieve's candidate list happened to include, so quantum
found a slightly cheaper feasible placement than the classical run's
reported best for this scenario --- not a discrepancy in either solver's
dispatch model.}
\end{table}
 in the Phase 3 report.
% Requires: booktabs, float (already loaded in the Phase 2 main doc)

\begin{table}[htbp]
\centering
\small
\begin{tabular}{llllll}
\toprule
\textbf{Scenario} & \textbf{Method} & \textbf{Siting Buses} & \textbf{24h Cost (\$)} & \textbf{Runtime} & \textbf{Gap vs.\ Classical} \\
\midrule
\multirow{3}{*}{No DC}
  & Classical siting (Benders/SCIP) & (1, 2, 3, 8)   & 97,229  & 16.3s  & (reference) \\
  & Quantum VQA --- Aer statevector ($n=20$) & (8, 11, 13, 14) & 97,229  & 102.4s & 0.000\% \\
  & QPU (IonQ Forte-1)               & ---            & ---     & ---    & --- \\
\addlinespace
\multirow{3}{*}{DC @ bus 4}
  & Classical siting (Benders/SCIP) & (3, 8, 11, 12) & 193,229 & 18.1s  & (reference) \\
  & Quantum VQA --- Aer statevector ($n=20$) & (4, 10, 11, 12) & 193,229 & 139.4s & 0.000\% \\
  & QPU (IonQ Forte-1)               & ---            & ---     & ---    & --- \\
\addlinespace
\multirow{3}{*}{DC @ 4, Gen2 outage}
  & Classical siting (Benders/SCIP) & (1, 2, 6, 7)   & 227,343 & 18.0s  & (reference) \\
  & Quantum VQA --- Aer statevector ($n=20$) & (4, 6, 9, 12) & 227,343 & 83.8s  & 0.000\% \\
  & QPU (IonQ Forte-1)               & ---            & ---     & ---    & --- \\
\addlinespace
\multirow{3}{*}{DC @ 4, heatwave}
  & Classical siting (Benders/SCIP) & (7, 9, 11, 12) & 203,675 & 34.6s  & (reference) \\
  & Quantum VQA --- Aer statevector ($n=20$) & (4, 9, 10, 14) & 203,404 & 76.8s  & $-$0.133\%\textsuperscript{\dag} \\
  & QPU (IonQ Forte-1)               & ---            & ---     & ---    & --- \\
\bottomrule
\end{tabular}
\caption{Cost and runtime comparison across four IEEE14 scenarios. Classical
siting uses Benders decomposition with SCIP subproblems (batch-Benders,
2-min time limit, early-stall detection). Quantum VQA uses a 19-qubit
butterfly ansatz (COBYLA, Aer CPU statevector simulator), evaluating the top
$n=20$ candidates from the quantum sieve classically (UC second stage). Gap
vs.\ Classical is the percentage cost difference relative to the classical
result; 0.000\% means the two methods matched to the nearest cent at this
problem size. QPU rows are placeholders (real Forte-1 hardware runs not yet
executed for this comparison --- budget-limited to a small number of real
QPU submissions). Gen2 outage (full 24h horizon) and the heatwave demand
scaling are independent scenarios, not combined. \textsuperscript{\dag}Negative
gap: the classical Benders search stopped early via its stall-detection
heuristic (no cost improvement for several iterations) before reaching the
placement the quantum sieve's candidate list happened to include, so quantum
found a slightly cheaper feasible placement than the classical run's
reported best for this scenario --- not a discrepancy in either solver's
dispatch model.}
\end{table}
